\documentclass[11pt,a4paper]{article}

\usepackage{fontspec}
\usepackage{polyglossia}
\usepackage{geometry}
\usepackage{graphicx}
\usepackage{xcolor}
\usepackage{booktabs}
\usepackage{array}
\usepackage{tabularx}
\usepackage{titlesec}
\usepackage{enumitem}
\usepackage{hyperref}
\usepackage{caption}
\usepackage{fancyhdr}
\usepackage{lastpage}

\setmainlanguage{english}
\setmainfont{Arial}

\geometry{top=20mm,bottom=20mm,left=18mm,right=18mm}
\setlength{\headheight}{14pt}
\graphicspath{{outputs/charts/}}

\definecolor{ink}{HTML}{17202A}
\definecolor{muted}{HTML}{5D6D7E}
\definecolor{blue}{HTML}{2D67B4}
\definecolor{green}{HTML}{238C5F}
\definecolor{red}{HTML}{C44348}
\definecolor{panel}{HTML}{F4F6F8}

\hypersetup{
    colorlinks=true,
    linkcolor=blue,
    urlcolor=blue
}

\pagestyle{fancy}
\fancyhf{}
\lhead{\textcolor{muted}{SuperStore Sales Performance Analysis}}
\rhead{\textcolor{muted}{2011--2014}}
\cfoot{\textcolor{muted}{\thepage}}
\renewcommand{\headrulewidth}{0pt}

\titleformat{\section}
  {\Large\bfseries\color{ink}}
  {}
  {0pt}
  {}

\titleformat{\subsection}
  {\large\bfseries\color{ink}}
  {}
  {0pt}
  {}

\setlist[itemize]{leftmargin=*, itemsep=2pt, topsep=4pt}
\setlength{\parindent}{0pt}
\setlength{\parskip}{6pt}
\captionsetup{font=small,labelfont=bf}

\newcommand{\metric}[2]{
    \begin{minipage}[t]{0.18\textwidth}
        \colorbox{panel}{
            \begin{minipage}[t][25mm][c]{0.92\linewidth}
                \centering
                {\Large\bfseries\color{ink} #1}\\[2mm]
                {\small\color{muted} #2}
            \end{minipage}
        }
    \end{minipage}
}

\begin{document}

\begin{titlepage}
    \thispagestyle{fancy}
    \vspace*{8mm}
    {\Huge\bfseries\color{ink} SuperStore Sales Performance Analysis\par}
    \vspace{3mm}
    {\Large\color{muted} Analytical report on sales, profit, and discount performance from 2011 to 2014\par}
    \vspace{12mm}

    \metric{\$12.6M}{Sales}
    \hfill
    \metric{\$1.47M}{Profit}
    \hfill
    \metric{11.6\%}{Profit margin}
    \hfill
    \metric{24.5\%}{Loss-making rows}
    \hfill
    \metric{51,290}{Rows}

    \vspace{14mm}
    \section*{Executive Summary}
    Global Superstore sales grew steadily from \$2.26M in 2011 to \$4.30M in 2014.
    Profit also increased, but overall profitability remains moderate at an 11.6\% profit margin.
    The main business risk is not demand volume; it is discount discipline. Nearly one quarter of transaction rows are loss-making, and discounts above 20\% systematically push orders into negative margin territory.

    \vspace{4mm}
    \begin{tabularx}{\textwidth}{>{\bfseries}p{42mm}X}
        \toprule
        Business question & Which markets, categories, and discount policies generate profit, and where does the company lose money? \\
        Analysis period & 2011--2014 \\
        Data scope & 25,035 orders, 795 customers, 7 markets \\
        Main risk & High discounts and concentrated losses in specific countries and sub-categories \\
        Key recommendation & Cap discounts above 20\% and review the Tables sub-category strategy \\
        \bottomrule
    \end{tabularx}

    \vfill
    {\color{muted} Prepared as a data analytics portfolio case study.}
\end{titlepage}
\setcounter{page}{2}

\section{1. Sales and Profit Trend}

Sales increased by 90.3\% between 2011 and 2014, from \$2.26M to \$4.30M.
Profit grew from \$249K to \$504K over the same period.
However, profit did not scale as efficiently as revenue, which suggests that discounting and loss-making categories limit operating leverage.

\begin{figure}[h]
    \centering
    \includegraphics[width=\textwidth]{sales_profit_trend.png}
    \caption{Sales and profit by year}
\end{figure}

\begin{tabularx}{\textwidth}{lrrrr}
    \toprule
    Year & Sales & Profit & Orders & Profit margin \\
    \midrule
    2011 & \$2,259,511 & \$248,941 & 4,440 & 11.0\% \\
    2012 & \$2,677,493 & \$307,415 & 5,343 & 11.5\% \\
    2013 & \$3,405,860 & \$408,513 & 6,721 & 12.0\% \\
    2014 & \$4,300,041 & \$504,166 & 8,531 & 11.7\% \\
    \bottomrule
\end{tabularx}

\section{2. Market Performance}

Profit is concentrated in APAC, EU, and the US. These markets provide the largest absolute profit contribution.
Canada has the strongest margin at 26.6\%, but its sales base is too small to act as the primary growth driver.
EMEA is the weakest market by profitability, generating only \$44K in profit at a 5.4\% margin.

\begin{figure}[h]
    \centering
    \includegraphics[width=\textwidth]{profit_by_market.png}
    \caption{Profit by market}
\end{figure}

\begin{tabularx}{\textwidth}{lrrrr}
    \toprule
    Market & Sales & Profit & Avg. discount & Profit margin \\
    \midrule
    APAC & \$3,585,833 & \$437,578 & 14.9\% & 12.2\% \\
    EU & \$2,938,139 & \$372,830 & 10.3\% & 12.7\% \\
    US & \$2,297,354 & \$286,397 & 15.6\% & 12.5\% \\
    LATAM & \$2,164,687 & \$221,643 & 13.6\% & 10.2\% \\
    Africa & \$783,776 & \$88,872 & 15.7\% & 11.3\% \\
    EMEA & \$806,184 & \$43,898 & 19.6\% & 5.4\% \\
    Canada & \$66,932 & \$17,817 & 0.0\% & 26.6\% \\
    \bottomrule
\end{tabularx}

\section{3. Category Performance}

Technology and Office Supplies perform similarly, with profit margins of 14.0\% and 13.7\%.
Furniture is materially weaker at a 7.0\% margin.
The main issue inside Furniture is the Tables sub-category, which generated a \$64K loss.

\begin{figure}[h]
    \centering
    \includegraphics[width=0.92\textwidth]{margin_by_category.png}
    \caption{Profit margin by category}
\end{figure}

\begin{tabularx}{\textwidth}{lrrrr}
    \toprule
    Category & Sales & Profit & Avg. discount & Profit margin \\
    \midrule
    Technology & \$4,744,691 & \$663,779 & 13.5\% & 14.0\% \\
    Office Supplies & \$3,787,330 & \$518,474 & 13.7\% & 13.7\% \\
    Furniture & \$4,110,884 & \$286,782 & 16.8\% & 7.0\% \\
    \bottomrule
\end{tabularx}

\section{4. Discount Impact}

The strongest finding is the relationship between discounts and profitability.
Orders with no discount have a 25.3\% profit margin.
At 20--30\% discounts, margin already turns negative at -5.5\%.
At discounts above 50\%, losses become severe, with a -111.0\% profit margin.

\begin{figure}[h]
    \centering
    \includegraphics[width=\textwidth]{margin_by_discount.png}
    \caption{Profit margin by discount bucket}
\end{figure}

\begin{tabularx}{\textwidth}{lrrrr}
    \toprule
    Discount bucket & Sales & Profit & Orders & Profit margin \\
    \midrule
    0\% & \$6,992,734 & \$1,770,695 & 15,211 & 25.3\% \\
    0--10\% & \$1,962,633 & \$339,767 & 3,028 & 17.3\% \\
    10--20\% & \$1,757,296 & \$173,255 & 4,363 & 9.9\% \\
    20--30\% & \$382,540 & -\$21,156 & 843 & -5.5\% \\
    30--50\% & \$1,176,078 & -\$380,945 & 3,674 & -32.4\% \\
    50\%+ & \$371,624 & -\$412,582 & 2,369 & -111.0\% \\
    \bottomrule
\end{tabularx}

\subsection{Discount impact by category}

The overall discount chart includes all markets, countries, categories, segments, and years.
The category-level view confirms that the discount pattern is broad-based: every major category becomes loss-making once discounts reach the 20--30\% range.
Furniture is the weakest category before heavy discounts as well, with only a 4.6\% margin at 10--20\% discounts.

\begin{figure}[h]
    \centering
    \includegraphics[width=\textwidth]{discount_impact_by_category.png}
    \caption{Profit margin by discount bucket and category}
\end{figure}

\subsection{Discount impact by market}

The market-level view shows that the same discount threshold is visible across most major markets.
APAC, EU, US, and LATAM are profitable up to the 10--20\% bucket, but each turns negative at 20--30\%.
Africa and EMEA have fewer intermediate discount buckets in the dataset, but their 50\%+ discount buckets show severe losses.
Canada has no discounted sales in this dataset, so only the 0\% bucket is available.

\begin{figure}[h]
    \centering
    \includegraphics[width=\textwidth]{discount_impact_by_market.png}
    \caption{Profit margin by discount bucket and market}
\end{figure}

\section{5. Loss Concentration by Country}

Losses are concentrated in countries with high average discounts.
Turkey, Nigeria, and the Netherlands are among the largest negative-profit countries.
This points to a commercial policy issue, not just a local demand or operations issue.

\begin{tabularx}{\textwidth}{lrrrr}
    \toprule
    Country & Sales & Profit & Avg. discount & Profit margin \\
    \midrule
    Turkey & \$108,524 & -\$98,447 & 60.0\% & -90.7\% \\
    Nigeria & \$54,347 & -\$80,751 & 70.0\% & -148.6\% \\
    Netherlands & \$77,520 & -\$41,070 & 48.2\% & -53.0\% \\
    Honduras & \$90,137 & -\$29,482 & 40.7\% & -32.7\% \\
    Pakistan & \$58,873 & -\$22,447 & 44.5\% & -38.1\% \\
    \bottomrule
\end{tabularx}

\section{Recommendations}

\begin{itemize}
    \item Cap discounts above 20\%, especially in countries and categories with negative margins.
    \item Review the Tables strategy, including pricing, discount rules, logistics, and product mix.
    \item Scale proven practices from APAC, EU, and the US because these markets generate the largest absolute profit.
    \item Audit EMEA, Turkey, and Nigeria at the pricing and promotional-policy level.
    \item Use Canada as a profitability benchmark, but not as the main growth driver because of its small sales base.
\end{itemize}

\section{Appendix: Reproducibility}

The analysis can be reproduced with:

\begin{verbatim}
python scripts/analyze_superstore.py
\end{verbatim}

\end{document}
